\section{Results}
\label{results}

This section first gives the overall picture, then explains why models have to be compared
split by split, and then goes through the classical models, the deep models, the stack and the
label-free models in turn.

\subsection{Overall comparison}

Table~\ref{tab:headline} gives the test metrics of every model as the mean and standard deviation
over the five split seeds, and Figure~\ref{fig:seeds} shows each seed separately. PR-AUC and
ROC-AUC measure the ranking and need no threshold. F1, precision and recall are measured at each
model's own validation-chosen threshold, so they depend on that choice as well as on the model.

\begin{table}[tb]
\centering
\footnotesize
\setlength{\tabcolsep}{4pt}
\begin{tabular}{lcccccr}
\toprule
\textbf{Model} & \textbf{PR-AUC} & \textbf{ROC-AUC} & \textbf{F1} & \textbf{Precision} & \textbf{Recall} & \textbf{s/seed} \\
\midrule
\tabinput{includes/tab_headline}
\bottomrule
\end{tabular}
\caption{Test metrics over five split seeds (mean $\pm$ standard deviation), sorted by PR-AUC. The
last column is the wall-clock time per seed, tuning included (classical models on a laptop CPU,
deep models on a GPU). The zero-shot time is for scoring all 17,599 ads once; the stack takes
under a second.}
\label{tab:headline}
\end{table}

\begin{figure}[tb]
\centering
\includegraphics[width=0.85\linewidth]{test_pr_auc_seeds.pdf}
\caption{Test PR-AUC of every model on every split seed. Dots are seeds, the black bar is the
mean and the dashed line is the fraud rate.}
\label{fig:seeds}
\end{figure}

Three groups stand out. The supervised text models (the stack, the logistic regressions, the MLP,
DistilBERT and Complement NB) reach 0.77 to 0.88. The frozen embeddings and the decision tree sit
in the middle, at 0.52 and 0.40. The three models that never learn from fraud labels, namely the
zero-shot language model, the one-class SVM and the dummy, stay at the 0.049 floor. The gap
between ROC-AUC and PR-AUC also confirms the choice of metric: the decision tree has a respectable
ROC-AUC of 0.843 but a PR-AUC of only 0.402.

\subsection{Comparing models across splits}

Within the top group, the standard deviations (0.03 to 0.08) are as large as most of the gaps
between the means, and in Figure~\ref{fig:seeds} the dots of the top five models overlap. Read
this way, the top models cannot be told apart. However, every model is evaluated on the same five
splits, and a split that is hard for one model is hard for all of them. Comparing two models seed
by seed cancels this shared difficulty, so Table~\ref{tab:paired} gives the per-seed differences
for the comparisons that matter.

\begin{table}[tb]
\centering
\small
\begin{tabular}{lrrrrrrc}
\toprule
\textbf{Pair (A $-$ B)} & \textbf{s0} & \textbf{s1} & \textbf{s2} & \textbf{s3} & \textbf{s4} & \textbf{Mean} & \textbf{A wins} \\
\midrule
\tabinput{includes/tab_paired}
\bottomrule
\end{tabular}
\caption{Paired differences in test PR-AUC per split seed.}
\label{tab:paired}
\end{table}

\subsection{Classical models}

The paired view separates models that the unpaired view could not. BM25 LR beats TF-IDF LR on all
five seeds (mean +0.026), although their means differ by less than one standard deviation, so
saturating repeated words and discounting long ads helps a little but consistently. The MLP is
level with BM25 LR (mean difference $-$0.001, 3 wins of 5) at about six times the training time,
so its extra capacity buys nothing here.

Complement NB loses to TF-IDF LR on 4 of 5 seeds (mean $-$0.048) but has the smallest spread of
any supervised model (standard deviation 0.029). It is also the most sensitive to its single
hyper-parameter: on the seed-0 validation rows it scores 0.427 untuned ($\alpha = 1$) and 0.867
tuned ($\alpha = 0.01$), which illustrates why every model is tuned. The decision tree reaches only
0.402 and has the largest gap between training and test PR-AUC (0.24): a single tree splits on a
handful of words and overfits them.

\subsection{Deep models}

Fine-tuned DistilBERT reaches 0.816. It is below BM25 LR on all five seeds (mean $-$0.032) and
level with TF-IDF LR (it wins on one seed of five), while needing a GPU. Its training curves
(Figure~\ref{fig:curves}, right) show that more training would not change this: validation PR-AUC
jumps from about 0.6 after the first epoch to about 0.8 after the second and is flat from the
third epoch on, with the best epoch at 3, 4 or 5 depending on the seed. The chosen thresholds vary
from 0.25 to 0.94 across seeds, because the fine-tuned model pushes most scores close to 0 or 1,
so a small shift in the scores moves the threshold a lot.

The frozen MiniLM embeddings reach only 0.524. A single 384-dimensional summary of what an ad
means loses the specific words, such as pay amounts or ``data entry'', that separate fraud in this
dataset. Fine-tuning the whole encoder recovers them (0.816), at a much higher cost.

\begin{figure}[tb]
\centering
\begin{minipage}{0.52\linewidth}
\centering
\includegraphics[width=\linewidth]{pr_curves_seed0.pdf}
\end{minipage}\hfill
\begin{minipage}{0.45\linewidth}
\centering
\includegraphics[width=\linewidth]{bert_epochs.pdf}
\end{minipage}
\caption{Left: test precision-recall curves of seed 0. Right: DistilBERT validation PR-AUC per
epoch, one thin line per seed and the mean as a thick line.}
\label{fig:curves}
\end{figure}

\subsection{Stacking}

Since DistilBERT on its own is weaker than BM25 LR, the question is whether it adds anything when
the two are combined. It does: the stack reaches 0.878, the best result of the project. It beats
BM25 LR on all five seeds (+0.022 to +0.036) and DistilBERT on all five (+0.051 to +0.087), and
its precision of 0.902 is the highest of any model, at a recall similar to BM25 LR's. In the
precision-recall curves (Figure~\ref{fig:curves}, left) the stack lies on or above both of its
inputs over most of the recall range. Measured per standard deviation of its input, BM25 LR's
weight in the stacker is 3 to 6 times DistilBERT's on every seed, so the stack is mostly BM25 LR,
corrected by DistilBERT where the two disagree.

\subsection{Label-free models}

Neither label-free detector does much better than the dummy. The one-class SVM scores 0.057 with a
ROC-AUC of 0.47, which is chance level. To find out whether the features or the anomaly assumption
is to blame, we trained a supervised logistic regression on the same SVD features: it reaches a
seed-0 validation PR-AUC of 0.685, against the one-class SVM's 0.066. The information is therefore
in the features, and what fails is the assumption that fraudulent ads are unusual.

The zero-shot Qwen2.5-1.5B model scores 0.078, with a ROC-AUC of 0.49. It answers ``fraudulent''
for 98\% of the seed-0 test ads, and its scores, computed in bf16, take only 47 distinct values
over 3,520 ads, so its ranking is both biased and coarse. Section~\ref{discussion} returns to why
both detectors fail.
